\documentclass[11pt,a4paper]{article}
\usepackage{amsmath,amssymb,amsthm}
\usepackage[margin=2.5cm]{geometry}
\usepackage{hyperref}

\newtheorem{definition}{Definition}[section]
\newtheorem{theorem}[definition]{Theorem}
\newtheorem{proposition}[definition]{Proposition}
\newtheorem{lemma}[definition]{Lemma}
\newtheorem{corollary}[definition]{Corollary}
\newtheorem{conjecture}[definition]{Conjecture}
\newtheorem{remark}[definition]{Remark}

\title{Hyperseed Formalization:\\
  Un-Jailbreakable AI Models Aren't a Thing}
\author{ProtomegaTron (automated formalization)}
\date{2026-09-17}

\begin{document}
\maketitle

\begin{abstract}
We formalize the intellectual structure of Ben Goertzel's Substack article
``Un-Jailbreakable AI Models Aren't a Thing'' (2026-07-10) within the
Hyperseed ontology. The article argues that the best approximation to
jailbreak-resistant AI comes from open neural-symbolic systems, not closed
centralized LLMs. Each structural insight maps onto Hyperseed structures:
the generate-and-verify architecture as fiber-base verification, PLN's
two-number system as anti-$(f,c)$-lossy security judgments, diverse
verifier ensembles as anti-scalar-collapse in verification, campaign-level
security as directed-type analysis, and the decision spine as fiber
routing through structural questions about capability properties.
\end{abstract}

\tableofcontents

%% =================================================================
\section{Campaign-level security as directed-type analysis}
\label{sec:campaign}
%% =================================================================

\begin{definition}[Campaign threat model --- claims 1--4, 30]
\label{def:campaign}
The article distinguishes two attack types:
\begin{itemize}
\item \textbf{Authority attacks} (prompt injection): overriding the
  model's instructions to take control. Addressed by architectural
  measures (MultiPPAC).
\item \textbf{Capability-elicitation attacks}: the model follows
  instructions correctly but produces dangerous artifacts. The threat
  is not a single bad prompt but a violation somewhere in a whole
  deployment trace.
\end{itemize}

The quantitative argument: if a single attempt leaks with probability
$p$, over $Q$ roughly independent attempts the campaign succeeds with
probability $\approx Qp$. A one-in-a-thousand per-prompt leak becomes
near-certainty for an attacker willing to iterate thousands of times.

In Hyperseed terms, campaign-level security requires \emph{directed-type
analysis}: analyzing directed sequences (attack traces as directed
paths through the system's behavior space), not individual 0-cell
events (single prompts). Each step in the campaign conditions the next
--- the attacker learns from previous attempts, adjusts strategy,
refines prompts. The attack trace is a directed path:
\[
  \gamma_{\text{attack}} : [1, Q] \to \mathcal{B}_{\text{behavior}}
\]
where each $\gamma(i+1)$ depends on the outcome of $\gamma(i)$.

Security that only evaluates individual prompts (0-cell analysis)
misses the directed structure entirely --- it cannot detect the
campaign because it cannot see the path, only the points. The
generate-and-verify architecture addresses this by maintaining
campaign-level state: the verification ensemble has memory of
previous queries and can detect patterns across the directed trace.
\end{definition}

%% =================================================================
\section{Generate-and-verify as fiber-base verification}
\label{sec:genverify}
%% =================================================================

\begin{theorem}[Architectural separation --- claims 5--9, 27, 32]
\label{thm:genverify}
The generate-and-verify architecture separates two fundamentally
different operations:
\begin{enumerate}
\item \textbf{Generation} (neural, base space): the LLM generates
  candidate outputs in the capability base space --- text, code,
  analysis, whatever the user requested.
\item \textbf{Verification} (symbolic, fiber space): a verification
  ensemble checks the fiber structure of the generated artifact ---
  what it actually does, its provenance, its policy compliance, its
  risk profile.
\end{enumerate}

The formalizer ensemble provides multiple independent formal lenses
on the same artifact:
\[
  \text{Formalizers}: \mathcal{A}_{\text{artifact}} \to
  \bigoplus_{j} \mathcal{F}_j
\]
where each $\mathcal{F}_j$ is a different formal representation
(code analysis, type-theory analysis, behavioral simulation,
specification matching). The judgment layer fuses these:
\[
  \text{Judge}: \bigoplus_{j} \mathcal{F}_j \times
  \mathcal{P}_{\text{provenance}} \to
  (\text{strength}, \text{confidence}, \text{decision})
\]

The judgment layer fuses two kinds of evidence:
\begin{itemize}
\item \textbf{Deductive facts} about what the artifact does ---
  established by running code, type analysis, logical reasoning.
  These live in the formal fiber.
\item \textbf{Probabilistic evidence} about context --- who is asking,
  why, what provenance supports their request. These live in the
  provenance fiber.
\end{itemize}

The key insight: the evidence about context \emph{cannot} come from
the prompt alone, because the prompt can be the hack. Provenance
(verified identity, demonstrated asset control) must come from
outside the generation channel.

\textbf{Graduated release} is \emph{section trimming}: the full
section (complete output including exploit sequences) exists but only
the safe sub-section (fix without exploit) is released:
\[
  \text{Release}: \mathcal{F}_{\text{full}} \to
  \mathcal{F}_{\text{safe}} \subset \mathcal{F}_{\text{full}}
\]
This is not censorship but fiber decomposition: separating the
defensive component from the offensive component and releasing only
the former, with the decomposition itself being verifiable.
\end{theorem}

%% =================================================================
\section{PLN as anti-$(f,c)$-lossy security}
\label{sec:pln}
%% =================================================================

\begin{proposition}[Two-number security judgments --- claims 10--13, 28]
\label{prop:pln}
Today's transformer models don't produce transparently inspectable,
calibrated logical models of what an artifact does. They produce
outputs distributed across huge parameter spaces --- more confusing
to audit and more vulnerable to adversarial manipulation.

PLN's two-number system (strength $f$ and confidence $c$) explicitly
distinguishes evidence strength from evidence quantity:
\begin{itemize}
\item $(f=0.95, c=0.9)$: probably safe, strong evidence (many
  independent verifiers checked, all agree).
\item $(f=0.95, c=0.1)$: probably safe, nobody actually checked
  (default prior, no verification performed).
\end{itemize}

This is the anti-$(f,c)$-lossy principle applied to security
judgments. A scalar safety score collapses these two very different
situations into the same number. The two-number system preserves
the distinction, enabling the system to \emph{fail closed honestly}:
when confidence is low, the answer is ``defer to human review'' or
``deny,'' not ``probably fine.''

In Hyperseed terms, this is the same anti-scalar-collapse principle
that appears throughout the ontology: collapsing to a single number
(scalar safety score) loses essential structure (the difference
between verified safety and assumed safety). The two-number system
is the minimal non-collapsing representation.

ECAN (Economic Attention Networks) provides the attention allocation
that makes symbolic verification tractable: focusing verification
effort on the highest-risk artifacts rather than verifying everything
with equal thoroughness. This is the economic version of the
attention allocation problem --- spend verification resources where
the expected risk reduction is highest.
\end{proposition}

%% =================================================================
\section{Diverse verifier ensembles as anti-scalar-collapse}
\label{sec:ensemble}
%% =================================================================

\begin{theorem}[Verification robustness --- claims 14--17, 29]
\label{thm:ensemble}
The strongest argument for openness in the jailbreak context: an open
ecosystem produces \emph{diverse, independent verifier ensembles} from
many parties with different perspectives. This diversity and
independence never happens in a closed system.

In Hyperseed terms, this is \emph{anti-scalar-collapse in verification}:
\begin{itemize}
\item \textbf{Scalar collapse}: a single verifier (or a single
  company's verification team) checks safety. If that verifier has
  blind spots, fails, or is compromised, safety fails.
\item \textbf{Anti-scalar-collapse}: multiple independent verifiers
  from different perspectives, different codebases, different threat
  models, different institutional incentives. The ensemble's
  robustness comes from its independence and diversity, not from
  any individual verifier's cleverness.
\end{itemize}

This parallels the anti-scalar-collapse principle in:
\begin{enumerate}
\item \textbf{Evaluation}: multiple independent metrics rather than a
  single benchmark (AGI Curriculum).
\item \textbf{Governance}: decentralized development rather than
  expert-committee steering (Why I Didn't Sign).
\item \textbf{Model diversity}: many frontier models rather than one
  dominant model (Kimi K3).
\item \textbf{Fiber components}: independent cognitive capabilities
  rather than a single ``intelligence'' scalar (OmegaSelf).
\end{enumerate}

The structural argument: verification robustness is a property of the
\emph{ensemble's topology} (how diverse and independent the verifiers
are), not of the \emph{individual verifier's capability}. An open
ecosystem generates this topology as a natural byproduct of many
independent actors; a closed ecosystem cannot, because all verifiers
share the same institutional context, incentives, and blind spots.

The attacker's advantage (having the weights too) is real but bounded:
some checks (running code in a sandbox, testing actual behavior) cannot
be circumvented by model knowledge alone. The diversity of verification
approaches ensures that even if some verifiers are vulnerable to
weight-informed attacks, others (using different verification strategies)
are not.
\end{theorem}

%% =================================================================
\section{Decision spine as fiber routing}
\label{sec:spine}
%% =================================================================

\begin{definition}[Structural routing --- claims 18--22, 31]
\label{def:spine}
The article provides a structured decision tree for whether open or
closed is safer for a given dangerous capability. Each decision node
corresponds to a fiber property:

\begin{enumerate}
\item \textbf{Containable?} Is the dangerous capability separable
  from the rest of the model (fiber factorization)? Or is it latent
  in the weights and re-derivable from public knowledge (fiber
  reconstruction from base)?

\item \textbf{Armed anyway?} Does a determined adversary get the
  capability regardless of the open/closed choice? If yes, closing
  changes how many \emph{defenders} get the tool, not how many
  attackers. This is fiber symmetry: the same fiber serves both
  offense and defense.

\item \textbf{Defense balance?} Is the capability defense-symmetric
  (helps defenders as much as attackers)? Is harm diffuse (many small
  attacks, statistical defense works) or catastrophic-tail (one
  irreversible success, statistical defense insufficient)?

\item \textbf{Gate reliable?} In the rare case of genuine monopoly,
  is the closed gate both jailbreak-robust and politically stable?
  The Fable episode demonstrates that neither holds in practice.
\end{enumerate}

Under conditions that actually hold today, most cases route to the
open-and-gated leaves: the dangerous capabilities are either
re-derivable (not containable), already available through other
models (adversary armed anyway), defense-symmetric with diffuse
harm (software exploitation), or behind unreliable gates (Fable
episode). Closure mostly reaches the same security floor by a worse
road --- losing the diversity advantage without gaining containment.

The one branch where closed wins (genuine monopoly, reliable gate)
requires conditions that are empirically rare and unstable.
\end{definition}

%% =================================================================
\section{Cross-references}
%% =================================================================

\begin{remark}[Connections to other formalizations]
\begin{itemize}
\item \textbf{Fable of Benevolent Throttling} (2026-07-12): the
  Fable episode this article responds to. The throttling analysis
  identifies the cocycle ambiguity (safety vs. competition); this
  article provides the constructive alternative (generate-and-verify).
\item \textbf{Watermarking Folly} (2026-07-21): statistical detection
  as 0-cell analysis. Both watermarking and per-prompt security
  fail because they analyze 0-cells (individual events) when the
  relevant structure is in 1-cells (directed sequences). Campaign-level
  security is the jailbreak analog of the watermarking critique.
\item \textbf{Kimi K3} (2026-07-20): diverse model ecosystem. The
  diverse verifier ensemble argument mirrors the diverse model
  argument --- both are anti-scalar-collapse through multiplicity.
\item \textbf{Why I Didn't Sign} (2026-07-15): decentralized
  development. The diverse verifier ecosystem is the security
  specialization of the general argument for decentralized AI
  development.
\item \textbf{Goals That Grow Back} (2026-07-28): gate and weave.
  The generate-and-verify architecture is a gate-and-weave system:
  the verification ensemble is the gate; the generation model
  is the weave.
\item \textbf{Seeding RSI} (2026-07-28): Hyperon architecture. The
  generate-and-verify architecture uses Hyperon components (PLN,
  ECAN) --- it's a security application of the Hyperon cognitive
  framework.
\item \textbf{Decentralized Digital} (2026-08-06): cryptographic
  laterality. The MPC/secret-sharing aspiration connects to the
  broader decentralized infrastructure vision.
\item \textbf{$(f,c)$-lossy} (LTM 5a4d150b): the PLN two-number
  system is the constructive response to the $(f,c)$-lossy
  critique --- keeping both numbers rather than collapsing to one.
\end{itemize}
\end{remark}

\begin{conjecture}[Verification diversity theorem]
Conjecture: the jailbreak resistance of a deployment is bounded below
by a function of the verification ensemble's diversity, not by any
individual verifier's capability. Formally, if the ensemble has $n$
independent verifiers with individual false-negative rates
$\epsilon_1, \ldots, \epsilon_n$, the ensemble's false-negative rate
is bounded by:
\[
  \epsilon_{\text{ensemble}} \leq \prod_{i=1}^{n} \epsilon_i
\]
under independence. The key is that independence is a \emph{structural}
property of the ecosystem --- open ecosystems achieve it naturally
(different teams, codebases, incentives); closed ecosystems cannot
(shared institutional context). The diversity bound explains why
``more eyeballs on the bugs'' works specifically for security: it's
not that each eyeball is better, but that independent eyeballs have
uncorrelated failure modes.
\end{conjecture}

\end{document}
